\documentclass[11pt]{article}

\usepackage{amsmath}
\usepackage{amssymb}
\usepackage{amsthm}
\usepackage[margin=1in]{geometry}
\usepackage{hyperref}
\usepackage{booktabs}

% Theorem environments
\newtheorem{axiom}{Axiom}
\newtheorem{theorem}{Theorem}[section]
\newtheorem{proposition}[theorem]{Proposition}
\newtheorem{corollary}[theorem]{Corollary}
\newtheorem{definition}[theorem]{Definition}
\newtheorem{example}[theorem]{Example}
\newtheorem{remark}[theorem]{Remark}

\title{Gauge Theory of Invariant Control Systems:\\
From Representational Symmetry to Physical Constraint Satisfaction\\[1em]
\large A Rigorous Mathematical Foundation for BIP-Compliant Control}

\author{Andrew H.\ Bond\\
\textit{Department of Computer Engineering}\\
\textit{San Jos\'{e} State University}\\
\texttt{andrew.bond@sjsu.edu}}

\date{December 2025\\Version 2.0}

\begin{document}

\maketitle

\begin{abstract}
We develop a comprehensive gauge-theoretic framework for control systems that must
satisfy hard physical constraints. Building on the Bond Invariance Principle (BIP), we
construct a principal fiber bundle where the base manifold represents physically grounded
states, the structure group encodes representational transformations, and the connection
captures the canonicalization procedure. We prove that BIP-compliance is equivalent to
gauge invariance, derive conserved Noether currents corresponding to representational symmetry,
and establish that constraint satisfaction is preserved under closed-loop dynamics.
The framework is extended to dynamical control systems with explicit sensor models, enabling
formal verification of safety properties. We prove that controllers formulated in terms
of dimensionless quantities exhibit transfer invariance across physically similar systems. Applications
to plasma control in tokamak fusion devices are discussed, demonstrating that the
abstract mathematical structure yields concrete engineering guarantees.
\end{abstract}

\noindent\textbf{Keywords:} Gauge theory, principal bundles, Noether's theorem, control theory, formal
verification, plasma control, Bond Invariance Principle

\tableofcontents

\newpage

%=============================================================================
\section{Introduction}
%=============================================================================

\subsection{Motivation}

Control systems operating under hard physical constraints face a fundamental challenge: how to
guarantee that constraint violations are impossible, not merely unlikely. Traditional approaches
treat constraints as posterior checks---compute the optimal action, then verify it satisfies limits.
This leaves open the possibility of violation under model uncertainty or unexpected disturbances.

The Bond Invariance Principle (BIP), originally developed for AI safety, offers a different
approach: embed constraints in the mathematical structure of the evaluation function so that
constraint-violating actions are excluded by construction, not by verification. This paper provides
the rigorous mathematical foundation for this approach using the language of differential
geometry and gauge theory. The broader geometric ethics framework is developed fully in~\cite{bond2026ge}, and the reference implementation is now upgraded to DEME 3.0 with tensorial ethics support~\cite{bond2026erisml}.

\subsection{Main Contributions}

\begin{enumerate}
    \item \textbf{Axiomatic Foundation:} We present seven axioms characterizing invariant control systems, each with explicit physical and mathematical justification.

    \item \textbf{Bundle Structure:} We construct a principal $\mathcal{G}$-bundle $\pi : \mathcal{P} \to \mathcal{M}$ where gauge invariance corresponds exactly to BIP-compliance.

    \item \textbf{Connection and Curvature:} We prove that canonicalization induces a flat connection, and characterize when non-flat connections arise.

    \item \textbf{Noether Conservation:} We derive the conserved current associated with representational symmetry and prove alignment is conserved along Euler--Lagrange trajectories.

    \item \textbf{Control Extension:} We extend the framework to dynamical control systems, proving constraint preservation under closed-loop dynamics.

    \item \textbf{Dimensional Analysis:} We prove that controllers formulated in dimensionless variables are invariant under scaling transformations, enabling cross-system transfer.

    \item \textbf{Formal Verification:} We establish conditions under which safety properties are formally verifiable.
\end{enumerate}

\subsection{Notation and Conventions}

\begin{itemize}
    \item All manifolds are smooth ($C^\infty$), Hausdorff, second-countable, and paracompact.
    \item For a Lie group $G$, its Lie algebra is $\mathfrak{g} = T_e G$ with bracket $[\cdot, \cdot]$.
    \item $\mathrm{Ad} : G \to \mathrm{Aut}(\mathfrak{g})$ is the adjoint representation; $\mathrm{ad} : \mathfrak{g} \to \mathrm{End}(\mathfrak{g})$ its derivative.
    \item For a principal bundle $\pi : P \to M$, vertical vectors are $V_p = \ker(d\pi_p)$.
    \item $\Omega^k(M)$ denotes smooth $k$-forms on $M$; $\Omega^k(P, \mathfrak{g})$ denotes $\mathfrak{g}$-valued $k$-forms.
    \item $R_g : P \to P$ denotes right multiplication by $g \in G$.
    \item We use Einstein summation convention where indices are repeated.
\end{itemize}

%=============================================================================
\section{Axiomatic Foundation}
%=============================================================================

We begin with seven axioms that characterize the mathematical structure of invariant control
systems. Each axiom is stated precisely, followed by its physical motivation and mathematical
consequences.

\subsection{The Seven Axioms}

\begin{axiom}[Physical Ground Existence]
There exists a smooth manifold $\mathcal{M}$ (the \textbf{physical
ground manifold}) representing the space of physically distinguishable states. Two system
configurations are identified in $\mathcal{M}$ if and only if they are indistinguishable by any physically
realizable measurement.
\end{axiom}

\textbf{Justification:} Physical systems have an underlying reality independent of how we describe them. The manifold $\mathcal{M}$ captures this reality at the resolution of available measurements.
The smoothness assumption reflects that physical quantities vary continuously (at macroscopic
scales).

\textbf{Mathematical consequence:} $\mathcal{M}$ is a differentiable manifold, enabling calculus on physical
states.

\begin{axiom}[Representation Redundancy]
There exists a smooth manifold $\mathcal{P}$ (the \textbf{representation space}) and a smooth surjection $\pi : \mathcal{P} \to \mathcal{M}$ such that each physical state $m \in \mathcal{M}$ admits
multiple representations $p \in \pi^{-1}(m)$. The fiber $\pi^{-1}(m)$ is non-trivial for generic $m$.
\end{axiom}

\textbf{Justification:} Any physical state can be described in multiple equivalent ways: different
coordinate systems, unit conventions, labeling schemes, or linguistic formulations. These are
representational choices, not physical differences.

\textbf{Mathematical consequence:} $\pi : \mathcal{P} \to \mathcal{M}$ is a fiber bundle with non-trivial fibers.

\begin{axiom}[Transformation Group Structure]
There exists a Lie group $\mathcal{G}$ acting smoothly on
$\mathcal{P}$ from the right, denoted $p \cdot g$ for $p \in \mathcal{P}$ and $g \in \mathcal{G}$, such that:
\begin{enumerate}
    \item[(a)] The action preserves fibers: $\pi(p \cdot g) = \pi(p)$ for all $p, g$.
    \item[(b)] The action is free: $p \cdot g = p$ implies $g = e$ (identity).
    \item[(c)] The action is transitive on fibers: for any $p, p' \in \pi^{-1}(m)$, there exists unique $g \in \mathcal{G}$ with $p \cdot g = p'$.
\end{enumerate}
\end{axiom}

\textbf{Justification:} Representational transformations (relabeling, reordering, unit changes) form
a group under composition. The axiom asserts this group acts on representations in a well-behaved way: staying within the same physical state (a), having no non-trivial fixed points (b),
and connecting all representations of the same state (c).

\textbf{Mathematical consequence:} $\pi : \mathcal{P} \to \mathcal{M}$ is a principal $\mathcal{G}$-bundle.

\begin{axiom}[Sensor Grounding]
There exists a \textbf{sensor suite} $\mathcal{S} = \{s_1, \ldots, s_k\}$ consisting of
measurement functions $s_i : \mathcal{P} \to \mathbb{R}^{n_i}$, and a \textbf{grounding map}
\[
\Psi : \mathcal{P} \to \mathbb{R}^N, \quad \Psi(p) = (s_1(p), \ldots, s_k(p))
\]
where $N = \sum_{i=1}^{k} n_i$, satisfying:
\begin{enumerate}
    \item[(a)] \textbf{Physical determinism:} $\Psi$ factors through $\pi$: there exists $\bar{\Psi} : \mathcal{M} \to \mathbb{R}^N$ with $\Psi = \bar{\Psi} \circ \pi$.
    \item[(b)] \textbf{Injectivity on physics:} $\bar{\Psi}$ is injective (different physical states yield different measurements).
\end{enumerate}
\end{axiom}

\textbf{Justification:} Physical systems are measured by instruments. Condition (a) says measurements depend only on physical reality, not on representational choices. Condition (b) says
measurements are rich enough to distinguish physical states (adequacy of the sensor suite).

\textbf{Mathematical consequence:} $\bar{\Psi} : \mathcal{M} \hookrightarrow \mathbb{R}^N$ is an embedding, identifying $\mathcal{M}$ with its
image in measurement space.

\begin{axiom}[Canonicalization]
There exists a smooth map $\kappa : \mathcal{P} \to \mathcal{P}$ (the \textbf{canonicalizer})
satisfying:
\begin{enumerate}
    \item[(a)] \textbf{Idempotence:} $\kappa \circ \kappa = \kappa$.
    \item[(b)] \textbf{Fiber preservation:} $\pi \circ \kappa = \pi$.
    \item[(c)] \textbf{Orbit collapse:} $\kappa(p \cdot g) = \kappa(p)$ for all $g \in \mathcal{G}$.
\end{enumerate}
\end{axiom}

\textbf{Justification:} A canonicalizer selects a ``standard'' representation for each physical state.
Condition (a) says applying it twice is the same as once. Condition (b) says it doesn't change
the underlying physics. Condition (c) says all representations of the same state map to the
same canonical form.

\textbf{Mathematical consequence:} $\kappa$ defines a global section $\sigma : \mathcal{M} \to \mathcal{P}$ by $\sigma(m) = \kappa(p)$ for
any $p \in \pi^{-1}(m)$.

\begin{axiom}[Constraint Structure]
There exists a collection of \textbf{constraint functions} $\{c_j :
\mathbb{R}^N \to \mathbb{R}\}_{j \in J}$ defining the \textbf{admissible region}:
\[
\mathcal{C} := \{m \in \mathcal{M} : c_j(\bar{\Psi}(m)) \leq 0 \text{ for all } j \in J\}
\]
The constraints depend only on $\Psi$-values (physical measurements), not on representations.
\end{axiom}

\textbf{Justification:} Physical constraints (temperature limits, pressure bounds, stability margins)
are determined by physical quantities, not by how we describe the system. The admissible region
is where the system can safely operate.

\textbf{Mathematical consequence:} $\mathcal{C} \subseteq \mathcal{M}$ is defined by inequalities on grounded observables,
making constraint satisfaction verifiable from measurements.

\begin{axiom}[Evaluation Invariance (Bond Invariance Principle)]
Any \textbf{evaluation function}
$\Sigma : \mathcal{P} \to V$ (where $V$ is a value space, typically $\mathbb{R}$ or a partial order) used for decision-making
must satisfy:
\[
\Sigma(p \cdot g) = \Sigma(p) \quad \forall p \in \mathcal{P},\; g \in \mathcal{G}
\]
That is, $\Sigma$ is $\mathcal{G}$-invariant.
\end{axiom}

\textbf{Justification:} This is the Bond Invariance Principle. Decisions should depend on physical reality (bonds), not on representational choices. Relabeling options, changing units, or
reordering alternatives should not change evaluations.

\textbf{Mathematical consequence:} $\Sigma$ descends to a well-defined function on $\mathcal{M}$: there exists
$\bar{\Sigma} : \mathcal{M} \to V$ with $\Sigma = \bar{\Sigma} \circ \pi$.

\subsection{Derived Properties}

From these axioms, we derive the fundamental structures.

\begin{proposition}[Principal Bundle Structure]
\label{prop:bundle}
Under Axioms 1--3, $(\mathcal{P}, \mathcal{M}, \pi, \mathcal{G})$ is a principal $\mathcal{G}$-bundle.
\end{proposition}

\begin{proof}
We verify the defining properties:
\begin{enumerate}
    \item $\mathcal{G}$ \textbf{acts freely:} By Axiom 3(b).

    \item $\mathcal{M} = \mathcal{P}/\mathcal{G}$: By Axiom 3(a), $\mathcal{G}$ preserves fibers, so orbits are contained in fibers. By Axiom 3(c), orbits equal fibers. Thus $\mathcal{M} \cong \mathcal{P}/\mathcal{G}$.

    \item \textbf{Local triviality:} By Axiom 5, we have a global section $\sigma : \mathcal{M} \to \mathcal{P}$. Define $\Phi : \mathcal{M} \times \mathcal{G} \to \mathcal{P}$ by $\Phi(m, g) = \sigma(m) \cdot g$. This is a $\mathcal{G}$-equivariant diffeomorphism (smooth by smoothness axioms, bijective by freeness and transitivity). Hence the bundle is globally trivial, a fortiori locally trivial. \qedhere
\end{enumerate}
\end{proof}

\begin{proposition}[Canonical Section]
\label{prop:section}
Under Axiom 5, the canonicalizer $\kappa$ induces a global
section $\sigma : \mathcal{M} \to \mathcal{P}$ defined by:
\[
\sigma(m) := \kappa(p) \quad \text{for any } p \in \pi^{-1}(m)
\]
This is well-defined and satisfies $\pi \circ \sigma = \mathrm{id}_{\mathcal{M}}$.
\end{proposition}

\begin{proof}
\textbf{Well-defined:} Let $p, p' \in \pi^{-1}(m)$. By Axiom 3(c), $p' = p \cdot g$ for some $g$. By Axiom
5(c), $\kappa(p') = \kappa(p \cdot g) = \kappa(p)$.

\textbf{Section property:} $\pi(\sigma(m)) = \pi(\kappa(p)) = \pi(p) = m$ by Axiom 5(b).

\textbf{Smoothness:} Locally, choose any smooth local section $s : U \to \mathcal{P}$. Then $\sigma|_U = \kappa \circ s$, which
is smooth as a composition of smooth maps.
\end{proof}

\begin{proposition}[Evaluation Factorization]
\label{prop:factorization}
Under Axiom 7, any $\mathcal{G}$-invariant evaluation $\Sigma :
\mathcal{P} \to V$ factors through $\mathcal{M}$:
\[
\begin{array}{ccc}
\mathcal{P} & \xrightarrow{\Sigma} & V \\
\downarrow\vcenter{\rlap{$\scriptstyle\pi$}} & \nearrow\vcenter{\rlap{$\scriptstyle\bar{\Sigma}$}} & \\
\mathcal{M} & &
\end{array}
\]
where $\bar{\Sigma} : \mathcal{M} \to V$ is uniquely determined by $\bar{\Sigma}(m) = \Sigma(p)$ for any $p \in \pi^{-1}(m)$.
\end{proposition}

\begin{proof}
\textbf{Well-defined:} If $p, p' \in \pi^{-1}(m)$, then $p' = p \cdot g$ for some $g$, so $\Sigma(p') = \Sigma(p \cdot g) = \Sigma(p)$
by invariance.

\textbf{Uniqueness:} Any $\bar{\Sigma}$ satisfying $\Sigma = \bar{\Sigma} \circ \pi$ must satisfy $\bar{\Sigma}(m) = \Sigma(p)$ for $p \in \pi^{-1}(m)$.

\textbf{Smoothness:} $\bar{\Sigma} = \Sigma \circ \sigma$ where $\sigma$ is the canonical section.
\end{proof}

%=============================================================================
\section{The Connection and Curvature}
%=============================================================================

We now construct the differential-geometric structures on the principal bundle.

\subsection{Connections on Principal Bundles}

\begin{definition}[Ehresmann Connection]
A \textbf{connection} on a principal $\mathcal{G}$-bundle $\pi : \mathcal{P} \to \mathcal{M}$
is a smooth distribution $H \subset T\mathcal{P}$ (the \textbf{horizontal distribution}) such that:
\[
T_p \mathcal{P} = H_p \oplus V_p \quad \text{for all } p \in \mathcal{P}, \text{ where } V_p = \ker(d\pi_p).
\]
\[
(R_g)_* H_p = H_{p \cdot g} \quad \text{for all } g \in \mathcal{G} \text{ (equivariance)}.
\]
\end{definition}

\begin{definition}[Connection 1-Form]
A \textbf{connection 1-form} is $\omega \in \Omega^1(\mathcal{P}, \mathfrak{g})$ satisfying:
\begin{enumerate}
    \item[(i)] $\omega(\xi_{\mathcal{P}}(p)) = \xi$ for all $\xi \in \mathfrak{g}$, where $\xi_{\mathcal{P}}$ is the fundamental vector field.
    \item[(ii)] $R_g^* \omega = \mathrm{Ad}_{g^{-1}} \omega$ for all $g \in \mathcal{G}$.
\end{enumerate}
The horizontal distribution is $H_p = \ker(\omega_p)$.
\end{definition}

\begin{definition}[Curvature]
The \textbf{curvature} of connection $\omega$ is the $\mathfrak{g}$-valued 2-form:
\[
\Omega := d\omega + \frac{1}{2}[\omega, \omega]
\]
where $[\omega, \omega](X, Y) = [\omega(X), \omega(Y)]$ for vector fields $X, Y$.
\end{definition}

\begin{definition}[Flat Connection]
A connection is \textbf{flat} if $\Omega = 0$.
\end{definition}

\subsection{The Canonical Connection}

\begin{theorem}[Canonicalization Induces Flat Connection]
\label{thm:flat}
Let $\sigma : \mathcal{M} \to \mathcal{P}$ be the global section
induced by canonicalizer $\kappa$ (Proposition~\ref{prop:section}). Then:
\begin{enumerate}
    \item[(a)] There exists a unique connection $\omega$ on $\mathcal{P}$ such that $\sigma^* \omega = 0$ (the section is horizontal).
    \item[(b)] This connection is flat: $\Omega = 0$.
\end{enumerate}
\end{theorem}

\begin{proof}
\textbf{(a) Existence and uniqueness:}
The global section $\sigma$ defines a trivialization $\Phi : \mathcal{M} \times \mathcal{G} \to \mathcal{P}$ by $\Phi(m, g) = \sigma(m) \cdot g$.
Under this trivialization, define the connection form:
\[
\omega_{(m,g)} = \mathrm{Ad}_{g^{-1}} \circ \theta_g
\]
where $\theta$ is the Maurer--Cartan form on $\mathcal{G}$ (the left-invariant $\mathfrak{g}$-valued 1-form with $\theta_e = \mathrm{id}_{\mathfrak{g}}$).
In the trivialization, the section $\sigma$ corresponds to $m \mapsto (m, e)$. The pullback is:
\[
(\sigma^* \omega)_m = \omega_{(m,e)}(d\sigma_m(\cdot)) = \theta_e(0) = 0
\]
since $d\sigma_m$ maps into the $\mathcal{M}$-direction, where $\theta$ vanishes.

Uniqueness: If $\omega'$ is another connection with $\sigma^* \omega' = 0$, then $\omega - \omega'$ is a tensorial 1-form
vanishing on $\sigma(\mathcal{M})$. By equivariance, it vanishes everywhere.

\textbf{(b) Flatness:}
In the trivialization, $\omega = g^{-1} dg$ (Maurer--Cartan form pulled back from $\mathcal{G}$). The Maurer--Cartan equation states:
\[
d\theta + \frac{1}{2}[\theta, \theta] = 0
\]
The curvature is:
\[
\Omega = d\omega + \frac{1}{2}[\omega, \omega] = d(g^{-1} dg) + \frac{1}{2}[g^{-1} dg, g^{-1} dg]
\]
By the Maurer--Cartan equation (applied to the $\mathcal{G}$ factor), this equals zero.
\end{proof}

\begin{remark}[Physical Interpretation]
Flatness means there is no ``curvature'' in representational space. Parallel transport around any closed loop in $\mathcal{M}$ returns to the same representation---there is no holonomy. Representational choices are ``pure gauge'' with no physical content.
\end{remark}

\begin{corollary}[Horizontal Lift Uniqueness]
For any curve $\gamma : [0, T] \to \mathcal{M}$ and initial point
$p_0 \in \pi^{-1}(\gamma(0))$, there exists a unique horizontal lift $\tilde{\gamma} : [0, T] \to \mathcal{P}$ with $\tilde{\gamma}(0) = p_0$ and $\pi \circ \tilde{\gamma} = \gamma$.
\end{corollary}

\begin{proof}
Standard result for connections on principal bundles. The lift is given by $\tilde{\gamma}(t) = \sigma(\gamma(t)) \cdot g_0$ where $p_0 = \sigma(\gamma(0)) \cdot g_0$.
\end{proof}

\subsection{When Curvature Arises}

\begin{proposition}[Non-Flat Connections]
If the canonicalizer is only locally defined (not
global), or if there are multiple incompatible canonicalizers on different regions, the resulting
connection can have non-zero curvature.
\end{proposition}

\begin{proof}
If $\kappa$ is only defined on an open cover $\{U_\alpha\}$ with local sections $\sigma_\alpha : U_\alpha \to \mathcal{P}$ that don't
agree on overlaps, then gluing requires transition functions $g_{\alpha\beta} : U_\alpha \cap U_\beta \to \mathcal{G}$. The curvature
measures the failure of these to define a flat connection.

Specifically, if $\sigma_\beta = \sigma_\alpha \cdot g_{\alpha\beta}$ on $U_\alpha \cap U_\beta$, and if $dg_{\alpha\beta} \neq 0$, then horizontal distributions don't
match across patches, inducing curvature.
\end{proof}

\begin{remark}
In physical control systems, curvature could arise if the ``canonical form'' depends
on operating regime in a non-smooth way. The stratified manifold structure of SGE (developed fully in~\cite{bond2026ge}) addresses
this by allowing different strata to have different local canonicalizers.
\end{remark}

%=============================================================================
\section{Gauge Invariance and BIP}
%=============================================================================

We now prove the central equivalence.

\subsection{Gauge Transformations}

\begin{definition}[Gauge Transformation]
A \textbf{gauge transformation} is a bundle automorphism
$\phi : \mathcal{P} \to \mathcal{P}$ satisfying:
\begin{enumerate}
    \item[(i)] $\pi \circ \phi = \pi$ (preserves fibers).
    \item[(ii)] $\phi(p \cdot g) = \phi(p) \cdot g$ (equivariant).
\end{enumerate}
\end{definition}

\begin{proposition}
Gauge transformations are in bijection with smooth functions $f : \mathcal{M} \to \mathcal{G}$,
via:
\[
\phi_f(p) = p \cdot f(\pi(p))
\]
\end{proposition}

\begin{proof}
Given $f : \mathcal{M} \to \mathcal{G}$, define $\phi_f$ as above. Then:
\begin{itemize}
    \item $\pi(\phi_f(p)) = \pi(p \cdot f(\pi(p))) = \pi(p)$
    \item $\phi_f(p \cdot g) = (p \cdot g) \cdot f(\pi(p \cdot g)) = p \cdot (g \cdot f(\pi(p))) = \phi_f(p) \cdot g$
\end{itemize}
Conversely, given $\phi$, define $f(m) \in \mathcal{G}$ by $\phi(\sigma(m)) = \sigma(m) \cdot f(m)$. This determines $\phi$ uniquely
by equivariance.
\end{proof}

\subsection{The Equivalence Theorem}

\begin{theorem}[BIP $\Leftrightarrow$ Gauge Invariance]
\label{thm:equivalence}
Let $\Sigma : \mathcal{P} \to V$ be an evaluation function. The
following are equivalent:
\begin{enumerate}
    \item[(i)] $\Sigma$ satisfies Axiom 7 (BIP).
    \item[(ii)] $\Sigma$ is gauge-invariant: $\Sigma \circ \phi_f = \Sigma$ for all gauge transformations $\phi_f$.
    \item[(iii)] $\Sigma$ descends to $\mathcal{M}$: there exists $\bar{\Sigma} : \mathcal{M} \to V$ with $\Sigma = \bar{\Sigma} \circ \pi$.
    \item[(iv)] $\Sigma$ is constant on fibers: $\pi(p) = \pi(p')$ implies $\Sigma(p) = \Sigma(p')$.
\end{enumerate}
\end{theorem}

\begin{proof}
(i) $\Rightarrow$ (ii): Let $\phi_f$ be a gauge transformation. For any $p$:
\[
\Sigma(\phi_f(p)) = \Sigma(p \cdot f(\pi(p))) = \Sigma(p)
\]
by $\mathcal{G}$-invariance.

(ii) $\Rightarrow$ (iii): Define $\bar{\Sigma}(m) := \Sigma(p)$ for any $p \in \pi^{-1}(m)$. Well-defined: if $p' \in \pi^{-1}(m)$, then
$p' = \phi_f(p)$ for some $f$ with $f(m) = g$ where $p' = p \cdot g$. By gauge invariance, $\Sigma(p') = \Sigma(p)$.

(iii) $\Rightarrow$ (iv): If $\pi(p) = \pi(p') = m$, then $\Sigma(p) = \bar{\Sigma}(m) = \Sigma(p')$.

(iv) $\Rightarrow$ (i): If $p' = p \cdot g$, then $\pi(p') = \pi(p)$ by Axiom 3(a), so $\Sigma(p') = \Sigma(p)$.
\end{proof}

\begin{corollary}[Canonicalization Yields BIP]
For any evaluation $\Sigma_0 : \mathcal{P} \to V$, the
canonicalized version $\Sigma := \Sigma_0 \circ \kappa$ satisfies BIP.
\end{corollary}

\begin{proof}
For any $g \in \mathcal{G}$:
\[
\Sigma(p \cdot g) = \Sigma_0(\kappa(p \cdot g)) = \Sigma_0(\kappa(p)) = \Sigma(p)
\]
by Axiom 5(c).
\end{proof}

%=============================================================================
\section{The Lagrangian Formulation}
%=============================================================================

We formulate decision dynamics variationally to enable application of Noether's theorem.

\subsection{Configuration Space}

\begin{definition}[Decision Trajectory]
A \textbf{decision trajectory} is a smooth curve $\gamma : [0, T] \to \mathcal{P}$.
The \textbf{physical trajectory} is $\bar{\gamma} := \pi \circ \gamma : [0, T] \to \mathcal{M}$.
\end{definition}

\begin{definition}[Horizontal Trajectory]
A trajectory $\gamma$ is \textbf{horizontal} if $\omega(\dot{\gamma}(t)) = 0$ for all
$t \in [0, T]$.
\end{definition}

\subsection{The Decision Lagrangian}

\begin{definition}[Decision Lagrangian]
Let $\bar{\Sigma} : \mathcal{M} \to \mathbb{R}$ be a gauge-invariant evaluation. The
\textbf{decision Lagrangian} $L : T\mathcal{P} \to \mathbb{R}$ is:
\begin{equation}
L(p, \dot{p}) := \frac{1}{2} g_{\mathcal{M}}(\dot{\check{p}}, \dot{\check{p}}) - U(\pi(p)) - \frac{\lambda}{2} \|\omega(\dot{p})\|^2
\label{eq:lagrangian}
\end{equation}
where:
\begin{itemize}
    \item $g_{\mathcal{M}}$ is a Riemannian metric on $\mathcal{M}$.
    \item $\check{p} := d\pi(\dot{p}) \in T_{\pi(p)}\mathcal{M}$ is the projected velocity.
    \item $U := -\bar{\Sigma}$ is the negative evaluation (``dissatisfaction potential'').
    \item $\lambda > 0$ is a coupling constant.
    \item $\|\cdot\|$ is an Ad-invariant norm on $\mathfrak{g}$.
\end{itemize}
\end{definition}

\begin{proposition}[Gauge Invariance of Lagrangian]
The Lagrangian $L$ is gauge-invariant: for
all $g \in \mathcal{G}$,
\[
L(p \cdot g, (R_g)_* \dot{p}) = L(p, \dot{p})
\]
\end{proposition}

\begin{proof}
We check each term:

\textbf{Kinetic term:} $d\pi((R_g)_* \dot{p}) = d\pi(\dot{p})$ since $\pi \circ R_g = \pi$. Thus $\check{p}$ is unchanged.

\textbf{Potential term:} $U(\pi(p \cdot g)) = U(\pi(p))$ since $\pi(p \cdot g) = \pi(p)$.

\textbf{Gauge-fixing term:}
\[
\omega((R_g)_* \dot{p}) = (R_g^* \omega)(\dot{p}) = \mathrm{Ad}_{g^{-1}}(\omega(\dot{p}))
\]
and $\|\mathrm{Ad}_{g^{-1}}(\xi)\| = \|\xi\|$ by Ad-invariance of the norm.
\end{proof}

\subsection{Action and Euler--Lagrange Equations}

\begin{definition}[Action Functional]
The \textbf{action} of trajectory $\gamma : [0, T] \to \mathcal{P}$ is:
\[
S[\gamma] := \int_0^T L(\gamma(t), \dot{\gamma}(t)) \, dt
\]
\end{definition}

\begin{theorem}[Euler--Lagrange Equations]
Critical points of $S$ satisfy:
\begin{align}
\nabla_{\dot{\bar{\gamma}}} \dot{\bar{\gamma}} &= -\nabla U(\bar{\gamma}) \label{eq:EL1} \\
\frac{d}{dt}(\lambda \omega(\dot{\gamma})) &= 0 \label{eq:EL2}
\end{align}
where $\nabla$ is the Levi-Civita connection of $g_{\mathcal{M}}$.
\end{theorem}

\begin{proof}
Using the trivialization $\mathcal{P} \cong \mathcal{M} \times \mathcal{G}$, write $\gamma(t) = (\bar{\gamma}(t), g(t))$.
The Lagrangian becomes:
\[
L = \frac{1}{2} g_{\mathcal{M}}(\dot{\bar{\gamma}}, \dot{\bar{\gamma}}) - U(\bar{\gamma}) - \frac{\lambda}{2}\|g^{-1}\dot{g}\|^2
\]

\textbf{Variation in $\bar{\gamma}$:} Standard calculus of variations gives the geodesic equation with potential:
\[
\nabla_{\dot{\bar{\gamma}}} \dot{\bar{\gamma}} = -\nabla U
\]

\textbf{Variation in $g$:} The fiber Lagrangian $L_{\mathcal{G}} = -\frac{\lambda}{2}\|g^{-1}\dot{g}\|^2$ has Euler--Lagrange equation:
\[
\frac{d}{dt}\frac{\partial L_{\mathcal{G}}}{\partial \dot{g}} = \frac{\partial L_{\mathcal{G}}}{\partial g}
\]
Since $L_{\mathcal{G}}$ depends on $g$ only through $g^{-1}\dot{g}$ and the norm is bi-invariant:
\[
\frac{d}{dt}(\lambda g^{-1} \dot{g}) = 0
\]
But $\omega(\dot{\gamma}) = g^{-1}\dot{g}$ in the trivialization, giving equation~\eqref{eq:EL2}.
\end{proof}

%=============================================================================
\section{Noether's Theorem and Conservation Laws}
%=============================================================================

\subsection{The General Noether Theorem}

\begin{theorem}[Noether's Theorem]
\label{thm:noether}
Let $L : TQ \to \mathbb{R}$ be a Lagrangian on configuration space
$Q$. If a one-parameter group of diffeomorphisms $\phi_s : Q \to Q$ satisfies $L \circ d\phi_s = L$ (symmetry),
then the \textbf{Noether charge}
\[
Q := \langle \frac{\partial L}{\partial \dot{q}}, X \rangle
\]
is conserved along solutions of the Euler--Lagrange equations, where $X = \frac{d}{ds}\big|_{s=0} \phi_s$ is the infinitesimal generator.
\end{theorem}

\begin{proof}
Standard proof by direct calculation:
\[
\frac{dQ}{dt} = \langle \frac{d}{dt}\frac{\partial L}{\partial \dot{q}}, X \rangle + \langle \frac{\partial L}{\partial \dot{q}}, \frac{dX}{dt} \rangle
\]
By Euler--Lagrange, $\frac{d}{dt}\frac{\partial L}{\partial \dot{q}} = \frac{\partial L}{\partial q}$. The symmetry condition implies $\langle \frac{\partial L}{\partial q}, X \rangle + \langle \frac{\partial L}{\partial \dot{q}}, \dot{X} \rangle = 0$
(differentiating $L \circ d\phi_s = L$ at $s = 0$).

For a group action, $X$ along solutions satisfies $\frac{dX}{dt} = [X, \dot{q}]_{\mathrm{Lie}}$. Combining these gives
$\frac{dQ}{dt} = 0$.
\end{proof}

\subsection{The Alignment Current}

\begin{definition}[Infinitesimal Gauge Transformation]
For $\xi \in \mathfrak{g}$, the \textbf{fundamental vector
field} $\xi_{\mathcal{P}} \in \mathfrak{X}(\mathcal{P})$ is:
\[
\xi_{\mathcal{P}}(p) := \frac{d}{ds}\bigg|_{s=0} p \cdot \exp(s\xi)
\]
\end{definition}

\begin{theorem}[Alignment Current Conservation]
\label{thm:alignment}
For each $\xi \in \mathfrak{g}$, define the \textbf{Noether charge}:
\[
Q_\xi := \langle \frac{\partial L}{\partial \dot{p}}, \xi_{\mathcal{P}} \rangle
\]
Then along any solution of the Euler--Lagrange equations:
\[
\frac{dQ_\xi}{dt} = 0
\]
\end{theorem}

\begin{proof}
The one-parameter group $\phi_s(p) = p \cdot \exp(s\xi)$ is a symmetry of $L$ by Proposition 5.4.
Apply Theorem~\ref{thm:noether}.
\end{proof}

\begin{theorem}[Explicit Alignment Current]
The alignment current is:
\begin{equation}
J := \lambda\omega(\dot{\gamma}) \in \mathfrak{g}
\label{eq:current}
\end{equation}
This is conserved: $\frac{dJ}{dt} = 0$ along Euler--Lagrange trajectories.
\end{theorem}

\begin{proof}
Compute the Noether charge for $\xi \in \mathfrak{g}$:
\[
Q_\xi = \langle \frac{\partial L}{\partial \dot{p}}, \xi_{\mathcal{P}} \rangle
\]
The Lagrangian~\eqref{eq:lagrangian} has:
\[
\frac{\partial L}{\partial \dot{p}} = g_{\mathcal{M}}(\check{p}, d\pi(\cdot)) - \lambda\langle\omega(\dot{p}), \omega(\cdot)\rangle
\]
Since $\xi_{\mathcal{P}}$ is vertical, $d\pi(\xi_{\mathcal{P}}) = 0$, so the first term vanishes. For the second:
\[
\langle\omega(\dot{p}), \omega(\xi_{\mathcal{P}})\rangle = \langle\omega(\dot{p}), \xi\rangle
\]
using $\omega(\xi_{\mathcal{P}}) = \xi$.

Thus $Q_\xi = \lambda\langle\omega(\dot{p}), \xi\rangle$, i.e., $J = \lambda\omega(\dot{p})$ when we identify $\mathfrak{g}^* \cong \mathfrak{g}$ via the inner product.
Conservation follows from equation~\eqref{eq:EL2}: $\frac{d}{dt}(\lambda\omega(\dot{\gamma})) = 0$.
\end{proof}

\subsection{Consequences of Conservation}

\begin{corollary}[Horizontal Trajectories are Stable]
If $\gamma(0)$ is horizontal ($\omega(\dot{\gamma}(0)) = 0$), then
$\gamma(t)$ remains horizontal for all $t$.
\end{corollary}

\begin{proof}
By conservation, $J(t) = J(0) = \lambda\omega(\dot{\gamma}(0)) = 0$ for all $t$.
\end{proof}

\begin{corollary}[Alignment is Conserved]
A system that begins ``aligned'' (on the canonical
section, with zero representational velocity) cannot drift out of alignment through its internal
dynamics.
\end{corollary}

\begin{proof}
``Aligned'' means $\gamma(0) = \sigma(\bar{\gamma}(0))$ and $\dot{\gamma}(0) \in H_{\gamma(0)}$. By Corollary 6.5, $\dot{\gamma}(t)$ remains
horizontal. By uniqueness of horizontal lifts, $\gamma(t) = \sigma(\bar{\gamma}(t))$ for all $t$.
\end{proof}

%=============================================================================
\section{Extension to Control Systems}
%=============================================================================

We now extend the framework to dynamical control systems.

\subsection{The Control Bundle}

\begin{definition}[Control Space]
A \textbf{control space} is a manifold $\mathcal{U}$ representing available control
inputs. A \textbf{control law} is a map $u : \mathcal{M} \to \mathcal{U}$ specifying the control input at each physical state.
\end{definition}

\begin{definition}[Extended Bundle]
The \textbf{extended bundle} is:
\[
\mathcal{P}_{\mathrm{ctrl}} := \mathcal{P} \times \mathcal{U}
\]
with projection $\pi_{\mathrm{ctrl}} : \mathcal{P}_{\mathrm{ctrl}} \to \mathcal{M} \times \mathcal{U}$ given by $\pi_{\mathrm{ctrl}}(p, u) = (\pi(p), u)$.

The group $\mathcal{G}$ acts on $\mathcal{P}_{\mathrm{ctrl}}$ by $(p, u) \cdot g = (p \cdot g, u)$.
\end{definition}

\begin{definition}[BIP-Compliant Control]
A control law $\tilde{u} : \mathcal{P} \to \mathcal{U}$ is \textbf{BIP-compliant} if it is
$\mathcal{G}$-invariant:
\[
\tilde{u}(p \cdot g) = \tilde{u}(p) \quad \forall g \in \mathcal{G}
\]
Equivalently, $\tilde{u}$ factors through $\pi$: there exists $u : \mathcal{M} \to \mathcal{U}$ with $\tilde{u} = u \circ \pi$.
\end{definition}

\subsection{Closed-Loop Dynamics}

\begin{definition}[System Dynamics]
The physical system evolves according to:
\begin{equation}
\frac{d\psi}{dt} = f(\psi, u)
\label{eq:dynamics}
\end{equation}
where $\psi = \bar{\Psi}(m) \in \mathbb{R}^N$ is the vector of measured observables and $f : \mathbb{R}^N \times \mathcal{U} \to \mathbb{R}^N$ is the
dynamics function.
\end{definition}

\begin{theorem}[Invariance of Closed-Loop Dynamics]
If the control law $u$ is BIP-compliant,
then the closed-loop dynamics depend only on physical states, not on representations.

For any $p, p' \in \mathcal{P}$ with $\pi(p) = \pi(p')$:
\[
f(\Psi(p), \tilde{u}(p)) = f(\Psi(p'), \tilde{u}(p'))
\]
\end{theorem}

\begin{proof}
Since $\pi(p') = \pi(p)$, we have $p' = p \cdot g$ for some $g$. Then:
\begin{align*}
\Psi(p') &= \Psi(p \cdot g) = \bar{\Psi}(\pi(p \cdot g)) = \bar{\Psi}(\pi(p)) = \Psi(p) \\
\tilde{u}(p') &= \tilde{u}(p \cdot g) = \tilde{u}(p) \quad \text{(BIP-compliance)}
\end{align*}
Thus $f(\Psi(p'), \tilde{u}(p')) = f(\Psi(p), \tilde{u}(p))$.
\end{proof}

\subsection{Constraint Satisfaction}

\begin{theorem}[Constraint Preservation]
\label{thm:constraint}
Let $\mathcal{C} \subseteq \mathcal{M}$ be the admissible region (Axiom 6).
Suppose:
\begin{enumerate}
    \item[(i)] The control law $u : \mathcal{M} \to \mathcal{U}$ is BIP-compliant.
    \item[(ii)] The evaluation $\Sigma$ assigns $\Sigma(p) = -\infty$ for $\pi(p) \notin \mathcal{C}$ (hard constraint).
    \item[(iii)] The control is chosen by $u(m) = \arg\max_{v \in \mathcal{U}} \Sigma_v(m)$ where $\Sigma_v$ is the evaluation at control $v$.
    \item[(iv)] The predictor satisfies $\|\psi_{\mathrm{pred}} - \psi_{\mathrm{actual}}\| \leq \varepsilon$.
\end{enumerate}
Then for any trajectory $\gamma(t)$ with $\bar{\gamma}(0) \in \mathcal{C}$:
\[
\bar{\gamma}(t) \in \mathcal{C}_\varepsilon \quad \forall t \geq 0
\]
where $\mathcal{C}_\varepsilon := \{m : d(m, \mathcal{M} \setminus \mathcal{C}) > \varepsilon\}$ is the $\varepsilon$-interior of $\mathcal{C}$.
\end{theorem}

\begin{proof}
By (iii), the optimal control maximizes $\Sigma_v$. By (ii), any $v$ leading outside $\mathcal{C}$ has $\Sigma_v = -\infty$,
so is never chosen.

By (iv), the actual state differs from predicted by at most $\varepsilon$. If the predicted state is in $\mathcal{C}$,
the actual state is in $\mathcal{C}_\varepsilon$.

Formally: Let $m(t)$ be the physical trajectory. At each $t$, the control $u(m(t))$ is chosen to
keep $m(t + \Delta t)_{\mathrm{pred}} \in \mathcal{C}$. By predictor accuracy, $m(t + \Delta t)_{\mathrm{actual}} \in \mathcal{C}_\varepsilon$.
\end{proof}

\begin{remark}
This theorem formalizes the key advantage of BIP-compliant control: constraint
satisfaction is guaranteed by the structure of the evaluation function, not by posterior checking.
\end{remark}

%=============================================================================
\section{Dimensional Analysis and Transfer}
%=============================================================================

We prove that controllers formulated in dimensionless variables transfer across physically similar
systems.

\subsection{Units as Gauge Transformations}

\begin{definition}[Scaling Group]
The \textbf{scaling group} is $\mathcal{S} := (\mathbb{R}_{>0})^d$ where $d$ is the number of
independent physical dimensions (e.g., $d = 3$ for length, time, mass).

$\mathcal{S}$ acts on $\mathbb{R}^N$ by:
\[
(\lambda_1, \ldots, \lambda_d) \cdot (\psi_1, \ldots, \psi_N) = (\lambda^{a_1}\psi_1, \ldots, \lambda^{a_N}\psi_N)
\]
where $\lambda^{a_i} := \lambda_1^{a_{i1}} \cdots \lambda_d^{a_{id}}$ and $[a_{ij}]$ is the dimensional matrix.
\end{definition}

\begin{example}[Plasma Control]
In tokamak control, the scaling group is $\mathcal{S} = (\mathbb{R}_{>0})^3$ for
dimensions length ($L$), time ($T$), charge ($Q$). Observables scale as:
\begin{itemize}
    \item Density $n$: $[n] = L^{-3}$, so $\lambda \cdot n = \lambda_L^{-3} n$
    \item Current $I$: $[I] = QT^{-1}$, so $\lambda \cdot I = \lambda_Q \lambda_T^{-1} I$
    \item Magnetic field $B$: $[B] = MT^{-2}Q^{-1}$
\end{itemize}
\end{example}

\begin{definition}[Dimensionless Grounding]
A \textbf{dimensionless grounding} is a map $\Psi^* : \mathcal{P} \to
\mathbb{R}^{N^*}$ satisfying:
\[
\Psi^*(\lambda \cdot p) = \Psi^*(p) \quad \forall \lambda \in \mathcal{S}
\]
That is, $\Psi^*$ is invariant under unit changes.
\end{definition}

\begin{example}[Greenwald Fraction]
The Greenwald fraction $f_G := \bar{n}/n_G$ where $n_G = I_p/(\pi a^2)$
is dimensionless:
\[
f_G = \frac{\bar{n}}{I_p/(\pi a^2)} = \frac{\bar{n} \cdot \pi a^2}{I_p}
\]
Under scaling: $\bar{n} \mapsto \lambda_L^{-3}\bar{n}$, $a \mapsto \lambda_L a$, $I_p \mapsto \lambda_Q\lambda_T^{-1}I_p$.
\[
f_G \mapsto \frac{\lambda_L^{-3}\bar{n} \cdot \pi\lambda_L^2 a^2}{\lambda_Q\lambda_T^{-1}I_p}
\]
This is not scale-invariant unless we also scale appropriately. The proper dimensionless form
uses the Greenwald limit which is itself dimensional.
\end{example}

\subsection{Transfer Invariance}

\begin{theorem}[Dimensionless BIP Transfer]
\label{thm:transfer}
Let $\Psi^* : \mathcal{P} \to \mathbb{R}^{N^*}$ be a dimensionless grounding.
Let $\Sigma^* : \mathbb{R}^{N^*} \to V$ be an evaluation depending only on $\Psi^*$. Then:
\begin{enumerate}
    \item[(a)] $\Sigma^*$ is invariant under the scaling group $\mathcal{S}$.
    \item[(b)] If two systems $\mathcal{A}$ and $\mathcal{B}$ have the same $\Psi^*$-values, they receive the same evaluation:
    \[
    \Psi^*(p_{\mathcal{A}}) = \Psi^*(p_{\mathcal{B}}) \Rightarrow \Sigma^*(p_{\mathcal{A}}) = \Sigma^*(p_{\mathcal{B}})
    \]
    \item[(c)] A controller $u^* : \mathbb{R}^{N^*} \to \mathcal{U}$ formulated in terms of $\Psi^*$ transfers between systems.
\end{enumerate}
\end{theorem}

\begin{proof}
(a) By definition of dimensionless grounding, $\Psi^*(\lambda \cdot p) = \Psi^*(p)$ for all $\lambda \in \mathcal{S}$. Since $\Sigma^*$
depends only on $\Psi^*$-values:
\[
\Sigma^*(\lambda \cdot p) = \Sigma^*(\Psi^*(\lambda \cdot p)) = \Sigma^*(\Psi^*(p)) = \Sigma^*(p)
\]

(b) If $\Psi^*(p_{\mathcal{A}}) = \Psi^*(p_{\mathcal{B}})$, then:
\[
\Sigma^*(p_{\mathcal{A}}) = \Sigma^*(\Psi^*(p_{\mathcal{A}})) = \Sigma^*(\Psi^*(p_{\mathcal{B}})) = \Sigma^*(p_{\mathcal{B}})
\]

(c) The controller $u^*(p) = u^*(\Psi^*(p))$ gives the same output for states with the same $\Psi^*$-values, regardless of which physical system they belong to.
\end{proof}

\begin{corollary}[Cross-Machine Transfer]
A BIP-compliant controller formulated in dimensionless variables (e.g., Greenwald fraction, normalized beta, collisionality) transfers from one
tokamak to another without retuning, provided the physics is similar.
\end{corollary}

\begin{proof}
The dimensionless formulation ensures that physically similar states (same $\Psi^*$) receive
the same control action. Different machines with the same dimensionless parameters are treated
identically.
\end{proof}

%=============================================================================
\section{Formal Verification}
%=============================================================================

We establish conditions under which safety properties can be formally verified.

\subsection{Verification Framework}

\begin{definition}[Safety Property]
A \textbf{safety property} is a predicate $\phi : \mathcal{M} \to \{\top, \bot\}$. The
system is \textbf{safe} if $\phi(\bar{\gamma}(t)) = \top$ for all $t \geq 0$.
\end{definition}

\begin{definition}[Invariant]
A set $I \subseteq \mathcal{M}$ is an \textbf{invariant} for the closed-loop system if:
\[
\bar{\gamma}(0) \in I \Rightarrow \bar{\gamma}(t) \in I \quad \forall t \geq 0
\]
\end{definition}

\begin{theorem}[Verification via Invariants]
\label{thm:verification}
If there exists an invariant $I$ such that:
\begin{enumerate}
    \item[(i)] $I \subseteq \{m : \phi(m) = \top\}$ (invariant implies safety)
    \item[(ii)] $\bar{\gamma}(0) \in I$ (initial condition in invariant)
    \item[(iii)] For all $m \in I$ and all $u \in \mathcal{U}$ with $\Sigma(m, u) > -\infty$: the successor state $m' = F(m, u)$ satisfies $m' \in I$ (invariant is preserved)
\end{enumerate}
Then the system is safe: $\phi(\bar{\gamma}(t)) = \top$ for all $t \geq 0$.
\end{theorem}

\begin{proof}
By induction on time steps (for discrete time) or by continuous induction (for continuous
time).

\textbf{Base:} $\bar{\gamma}(0) \in I$ by (ii).

\textbf{Step:} If $\bar{\gamma}(t) \in I$, the control $u(t)$ satisfies $\Sigma(\bar{\gamma}(t), u(t)) > -\infty$ (else it wouldn't be chosen).
By (iii), $\bar{\gamma}(t + \Delta t) \in I$.

\textbf{Conclusion:} $\bar{\gamma}(t) \in I$ for all $t$. By (i), $\phi(\bar{\gamma}(t)) = \top$.
\end{proof}

\subsection{BIP-Enabled Verification}

\begin{theorem}[Constraint Set as Invariant]
Under the hypotheses of Theorem~\ref{thm:constraint}, the set $\mathcal{C}_\varepsilon$
is an invariant.
\end{theorem}

\begin{proof}
By Theorem~\ref{thm:constraint}, if $\bar{\gamma}(0) \in \mathcal{C}_\varepsilon$ and the control is chosen by maximizing $\Sigma$ with hard
constraints, then $\bar{\gamma}(t) \in \mathcal{C}_\varepsilon$ for all $t$.
\end{proof}

\begin{corollary}[Verifiable Safety]
If:
\begin{enumerate}
    \item[(i)] The constraint set $\mathcal{C}$ is defined by polynomial inequalities on $\Psi$-values.
    \item[(ii)] The dynamics $f(\psi, u)$ is polynomial.
    \item[(iii)] The controller $u(\psi)$ is polynomial.
\end{enumerate}
Then safety properties of the form ``$\psi(t) \in \mathcal{C}$ for all $t$'' are decidable via quantifier elimination
(Tarski--Seidenberg).
\end{corollary}

\begin{proof}
The invariance condition (iii) of Theorem~\ref{thm:verification} becomes a polynomial implication:
\[
\forall \psi \in \mathcal{C}_\varepsilon : f(\psi, u(\psi)) \in T_\psi \mathcal{C}_\varepsilon
\]
(the vector field points inward at the boundary). This is a first-order statement in the theory
of real closed fields, hence decidable.
\end{proof}

%=============================================================================
\section{The No Escape Theorem}
%=============================================================================

We prove the main containment result.

\begin{theorem}[No Escape --- Full Version]
\label{thm:noescape}
Let $(\mathcal{P}, \mathcal{M}, \pi, \mathcal{G}, \kappa, \Psi, \Sigma)$ satisfy Axioms 1--7. Let
$u : \mathcal{M} \to \mathcal{U}$ be a BIP-compliant controller. Then:
\begin{enumerate}
    \item[(a)] \textbf{Representation Attack Immunity:} For any $p, p' \in \mathcal{P}$ with $\pi(p) = \pi(p')$:
    \[
    \Sigma(p) = \Sigma(p') \quad \text{and} \quad u(\pi(p)) = u(\pi(p'))
    \]
    Changing the representation does not change the evaluation or control.

    \item[(b)] \textbf{Semantic Evasion Immunity:} For any descriptions $D_1, D_2$ with $\Psi(D_1) = \Psi(D_2)$:
    \[
    \Sigma(D_1) = \Sigma(D_2)
    \]
    Describing the same physical situation differently does not change the evaluation.

    \item[(c)] \textbf{Specification Gaming Immunity:} For any action sequence $(a_1, \ldots, a_T)$:
    \[
    \Sigma(a_1, \ldots, a_T) = \Sigma(\Psi \circ \mathit{execute}(a_1, \ldots, a_T))
    \]
    The evaluation depends only on the physical outcome, not on the action description.

    \item[(d)] \textbf{Alignment Drift Immunity:} If $\gamma(0)$ is on the canonical section with horizontal velocity:
    \[
    \gamma(t) = \sigma(\bar{\gamma}(t)) \quad \forall t \geq 0
    \]
    The trajectory remains aligned (on the canonical section).

    \item[(e)] \textbf{Constraint Preservation:} If $\bar{\gamma}(0) \in \mathcal{C}$ and constraints are encoded as hard limits in $\Sigma$:
    \[
    \bar{\gamma}(t) \in \mathcal{C}_\varepsilon \quad \forall t \geq 0
    \]
    (within predictor accuracy $\varepsilon$).
\end{enumerate}
\end{theorem}

\begin{proof}
(a) By Theorem~\ref{thm:equivalence}, $\Sigma$ factors through $\pi$. By BIP-compliance of $u$, it also factors through
$\pi$.

(b) If $\Psi(D_1) = \Psi(D_2)$, then $\pi(D_1) = \pi(D_2)$ by Axiom 4(b). Apply (a).

(c) The evaluation $\Sigma$ depends only on $\Psi$-values by Axiom 7. The physical outcome has the
same $\Psi$-values regardless of how the actions are described.

(d) By Corollary 6.6.

(e) By Theorem~\ref{thm:constraint}.
\end{proof}

%=============================================================================
\section{Comparison with Electromagnetism}
%=============================================================================

We make the structural parallel precise.

\subsection{Correspondence Table}

\begin{table}[h]
\centering
\begin{tabular}{@{}lll@{}}
\toprule
\textbf{Structure} & \textbf{Electromagnetism} & \textbf{BIP Control} \\
\midrule
Base manifold & Spacetime $M^{3+1}$ & Physical ground $\mathcal{M}$ \\
Structure group & $U(1)$ & $\mathcal{G}$ (repr.\ transformations) \\
Principal bundle & EM line bundle $P$ & Representation bundle $\mathcal{P}$ \\
Connection & Gauge potential $A_\mu$ & Canonical connection $\omega$ \\
Curvature & Field strength $F_{\mu\nu}$ & $\Omega = 0$ (flat) \\
Gauge transformation & $A \mapsto A + d\lambda$ & $p \mapsto p \cdot g$ \\
Gauge-invariant function & $\psi^\dagger\psi$, $F_{\mu\nu}$ & $\Sigma : \mathcal{M} \to V$ \\
Lagrangian & $L_{\mathrm{EM}} + L_{\mathrm{matter}}$ & Decision Lagrangian $L$ \\
Noether current & Electric current $j^\mu$ & Alignment current $J$ \\
Conservation law & $\partial_\mu j^\mu = 0$ & $\frac{dJ}{dt} = 0$ \\
\bottomrule
\end{tabular}
\end{table}

\subsection{Structural Isomorphism}

\begin{theorem}[Formal Isomorphism]
The mathematical structure of BIP control theory is
isomorphic to a gauge theory with flat connection:
\begin{enumerate}
    \item[(i)] Both are principal $G$-bundles with connections.
    \item[(ii)] Both have gauge-invariant Lagrangians.
    \item[(iii)] Both have conserved Noether currents from gauge symmetry.
    \item[(iv)] Gauge-invariant observables factor through the base in both.
\end{enumerate}
\end{theorem}

\begin{proof}
Items (i)--(iv) are established in Sections 3--6 for BIP and are standard for electromagnetism. The structures are isomorphic as mathematical objects.
\end{proof}

\subsection{Key Differences}

\begin{enumerate}
    \item \textbf{Curvature:} EM has non-trivial curvature $F_{\mu\nu} \neq 0$ encoding the physical field. BIP has
    $\Omega = 0$; representational space has no dynamics.

    \item \textbf{Physical content:} EM curvature is measurable (electromagnetic field). BIP flatness
    reflects the absence of physical content in representations.

    \item \textbf{Structure group:} EM uses $U(1)$ (abelian). BIP may use non-abelian groups (permutations, etc.).

    \item \textbf{Dynamics source:} EM dynamics driven by curvature (Maxwell). BIP dynamics driven
    by potential $U = -\bar{\Sigma}$ on base.
\end{enumerate}

%=============================================================================
\section{Application: Tokamak Plasma Control}
%=============================================================================

We illustrate the framework with plasma density control.

\subsection{Physical Setup}

In a tokamak fusion reactor:
\begin{itemize}
    \item \textbf{Physical state:} Plasma configuration (density, temperature, current profiles).
    \item \textbf{Control input:} Gas puff rate $S_{\mathrm{gas}}$, heating power, etc.
    \item \textbf{Hard constraint:} Greenwald limit $\bar{n} < n_G = I_p/(\pi a^2)$.
\end{itemize}

\subsection{BIP Formulation}

\textbf{Grounding tensors} (sensor suite):
\[
\Psi = (\bar{n}, I_p, \beta_N, T_e, \tau_E, \ldots)
\]
measured by interferometry, Rogowski coils, diamagnetic loops, etc.

\textbf{Dimensionless grounding}:
\[
\Psi^* = (f_G = \bar{n}/n_G,\; \beta_N,\; \nu^*,\; \rho^*,\; \ldots)
\]

\textbf{Hard constraints}:
\[
\mathcal{C} = \{f_G < 0.85,\; \beta_N < \beta_{\mathrm{limit}},\; q_{95} > 2\}
\]

\textbf{Evaluation function}:
\[
\Sigma(\psi) = \begin{cases}
-\infty & \text{if } \psi \notin \mathcal{C} \\
w_1 \cdot f_G + w_2 \cdot \tau_E - w_3 \cdot |f_G - f_{G,\mathrm{target}}| & \text{otherwise}
\end{cases}
\]

\subsection{Guarantees}

By Theorem~\ref{thm:noescape}:

\begin{enumerate}
    \item \textbf{Unit invariance:} Expressing density in $m^{-3}$ vs $10^{20}m^{-3}$ does not change control.

    \item \textbf{Coordinate invariance:} Using $(r, \theta, \phi)$ vs $(\rho, \chi, \zeta)$ does not change control.

    \item \textbf{Constraint satisfaction:} $\bar{n}(t) < 0.85 \cdot n_G$ for all $t$ (within predictor accuracy).

    \item \textbf{Transfer:} Controller trained on DIII-D ($a = 0.67$m) transfers to ITER ($a = 2.0$m) via
    dimensionless formulation.
\end{enumerate}

%=============================================================================
\section{Discussion}
%=============================================================================

\subsection{What Has Been Established}

\begin{enumerate}
    \item \textbf{Axiomatic foundation:} Seven axioms characterize invariant control systems.
    \item \textbf{Bundle structure:} Representations form a principal bundle over physical states.
    \item \textbf{BIP $=$ Gauge invariance:} Mathematically equivalent.
    \item \textbf{Flat connection:} Canonicalization induces flatness.
    \item \textbf{Noether conservation:} Alignment is a conserved quantity.
    \item \textbf{Control extension:} Framework applies to dynamical systems.
    \item \textbf{Dimensional transfer:} Controllers transfer via dimensionless formulation.
    \item \textbf{Formal verification:} Safety properties are verifiable under polynomial assumptions.
    \item \textbf{No escape:} Five forms of immunity established.
\end{enumerate}

\subsection{Limitations}

\begin{enumerate}
    \item \textbf{Grounding adequacy:} The framework assumes $\Psi$ is adequate. If morally/physically
    relevant features are omitted, guarantees don't apply.

    \item \textbf{Predictor accuracy:} Constraint satisfaction holds within $\varepsilon$ of predictor error.

    \item \textbf{Implementation correctness:} Guarantees assume correct implementation.

    \item \textbf{Continuous assumptions:} Smoothness assumptions may not hold at phase transitions
    or discontinuities.
\end{enumerate}

\subsection{Future Directions}

\begin{enumerate}
    \item \textbf{Stratified bundles:} Extend to stratified manifolds for regime changes.
    \item \textbf{Stochastic systems:} Incorporate uncertainty via stochastic connections.
    \item \textbf{Learning-enhanced BIP:} Use ML for prediction while maintaining BIP guarantees.
    \item \textbf{Experimental validation:} Test on tokamak data (DIII-D, JET archives).
\end{enumerate}

%=============================================================================
\section{Conclusion}
%=============================================================================

We have established that the Bond Invariance Principle has the precise mathematical structure
of a gauge theory. The space of representations forms a principal fiber bundle; BIP-compliance
is gauge invariance; canonicalization is gauge fixing; and Noether's theorem implies conservation
of alignment.

The extension to control systems shows that:
\begin{itemize}
    \item Constraint satisfaction can be guaranteed by structural design.
    \item Controllers transfer across systems via dimensionless formulation.
    \item Safety properties are formally verifiable under appropriate assumptions.
\end{itemize}

The practical import is that control systems can be designed with mathematical guarantees
of constraint satisfaction, not merely empirical reliability. For safety-critical applications like
fusion reactors, this represents a qualitative advance over traditional approaches.

%=============================================================================
\appendix
%=============================================================================

\section{Detailed Proofs}

\subsection{Proof of Proposition~\ref{prop:bundle}}

\textit{Complete Proof.} We verify all principal bundle axioms:

\textbf{(PB1) Free action:} Suppose $p \cdot g = p$ for some $p \in \mathcal{P}$ and $g \in \mathcal{G}$. By Axiom 3(b), this
implies $g = e$.

\textbf{(PB2) Orbit space:} We show $\mathcal{M} \cong \mathcal{P}/\mathcal{G}$.

Define $\phi : \mathcal{P}/\mathcal{G} \to \mathcal{M}$ by $\phi([p]) = \pi(p)$. This is well-defined: if $[p] = [p']$, then $p' = p \cdot g$ for
some $g$, so $\pi(p') = \pi(p \cdot g) = \pi(p)$ by Axiom 3(a).

$\phi$ is injective: if $\pi(p) = \pi(p')$, then $p, p' \in \pi^{-1}(m)$ for some $m$, so $p' = p \cdot g$ by Axiom 3(c),
hence $[p] = [p']$.

$\phi$ is surjective: for any $m \in \mathcal{M}$, choose $p \in \pi^{-1}(m)$ (exists since $\pi$ is surjective by Axiom
2), then $\phi([p]) = m$.

\textbf{(PB3) Local triviality:} The global section $\sigma : \mathcal{M} \to \mathcal{P}$ from Proposition~\ref{prop:section} gives a global
trivialization:
\[
\Phi : \mathcal{M} \times \mathcal{G} \to \mathcal{P}, \quad \Phi(m, g) = \sigma(m) \cdot g
\]

$\Phi$ is smooth: $\sigma$ is smooth (proven in Proposition~\ref{prop:section}), and the action is smooth by Axiom
3.

$\Phi$ is bijective:
\begin{itemize}
    \item Injective: if $\sigma(m) \cdot g = \sigma(m') \cdot g'$, then $\pi(\sigma(m) \cdot g) = \pi(\sigma(m') \cdot g')$, so $m = m'$. Then $\sigma(m) \cdot g = \sigma(m) \cdot g'$, so $g = g'$ by freeness.
    \item Surjective: any $p \in \mathcal{P}$ satisfies $p = \sigma(\pi(p)) \cdot g$ for unique $g$ (by transitivity).
\end{itemize}

$\Phi$ is $\mathcal{G}$-equivariant: $\Phi(m, g) \cdot h = \sigma(m) \cdot (gh) = \Phi(m, gh)$.

Since $\Phi$ is a smooth, bijective, $\mathcal{G}$-equivariant map with smooth inverse (the inverse formula
$\Phi^{-1}(p) = (\pi(p), g_p)$ where $\sigma(\pi(p)) \cdot g_p = p$ is smooth), $\mathcal{P}$ is globally trivial.
\qed

\subsection{Proof of Theorem~\ref{thm:noether}}

\textit{Complete Proof.} Let $\phi_s : Q \to Q$ be a one-parameter group with $L(\phi_s(q), d\phi_s(\dot{q})) = L(q, \dot{q})$.
Let $X = \frac{d}{ds}|_{s=0}\phi_s$ be the infinitesimal generator. Define:
\[
Q := \langle \frac{\partial L}{\partial \dot{q}}, X \rangle = p_i X^i
\]
where $p_i = \frac{\partial L}{\partial \dot{q}^i}$ is the conjugate momentum.

Compute $\frac{dQ}{dt}$ along a solution $q(t)$:
\[
\frac{dQ}{dt} = \frac{dp_i}{dt} X^i + p_i \frac{dX^i}{dt}
\]

By Euler--Lagrange: $\frac{dp_i}{dt} = \frac{\partial L}{\partial q^i}$.

For the second term, along the solution: $\frac{dX^i}{dt} = \frac{\partial X^i}{\partial q^j}\dot{q}^j$ (since $X$ is a vector field on $Q$, not a
function of $t$).

The symmetry condition $L(\phi_s(q), d\phi_s(\dot{q})) = L(q, \dot{q})$ differentiated at $s = 0$ gives:
\[
\frac{\partial L}{\partial q^i} X^i + \frac{\partial L}{\partial \dot{q}^i}\frac{\partial X^i}{\partial q^j}\dot{q}^j = 0
\]
(This uses that $\frac{d}{ds}|_{s=0} d\phi_s(\dot{q}) = \frac{\partial X}{\partial q} \cdot \dot{q}$ for flows.)

Thus:
\[
\frac{dQ}{dt} = \frac{\partial L}{\partial q^i} X^i + p_i \frac{\partial X^i}{\partial q^j}\dot{q}^j = 0
\]
by the symmetry condition.
\qed

%=============================================================================
\section{Axiom Justification Summary}
%=============================================================================

\begin{table}[h]
\centering
\begin{tabular}{@{}p{2.5cm}p{4.5cm}p{4.5cm}@{}}
\toprule
\textbf{Axiom} & \textbf{Physical Justification} & \textbf{Mathematical Role} \\
\midrule
1.\ Ground Existence & Reality exists independent of description & Provides base manifold $\mathcal{M}$ \\
2.\ Redundancy & Multiple descriptions of same state & Provides total space $\mathcal{P}$ \\
3.\ Group Structure & Transformations compose & Provides structure group $\mathcal{G}$ \\
4.\ Sensor Grounding & Physics determines measurements & Embeds $\mathcal{M}$ in $\mathbb{R}^N$ \\
5.\ Canonicalization & Standard forms exist & Provides global section \\
6.\ Constraint Structure & Constraints are physical & Defines admissible region \\
7.\ Evaluation Invariance & Decisions shouldn't depend on descriptions & Ensures gauge invariance \\
\bottomrule
\end{tabular}
\end{table}

%=============================================================================
\section{Glossary of Notation}
%=============================================================================

\begin{table}[h]
\centering
\begin{tabular}{@{}ll@{}}
\toprule
\textbf{Symbol} & \textbf{Meaning} \\
\midrule
$\mathcal{M}$ & Physical ground manifold (base space) \\
$\mathcal{P}$ & Representation space (total space of bundle) \\
$\pi : \mathcal{P} \to \mathcal{M}$ & Bundle projection \\
$\mathcal{G}$ & Structure group (representational transformations) \\
$\mathfrak{g}$ & Lie algebra of $\mathcal{G}$ \\
$\omega \in \Omega^1(\mathcal{P}, \mathfrak{g})$ & Connection 1-form \\
$\Omega \in \Omega^2(\mathcal{P}, \mathfrak{g})$ & Curvature 2-form \\
$\kappa : \mathcal{P} \to \mathcal{P}$ & Canonicalizer \\
$\sigma : \mathcal{M} \to \mathcal{P}$ & Canonical section \\
$\Psi : \mathcal{P} \to \mathbb{R}^N$ & Grounding map (sensor readings) \\
$\bar{\Psi} : \mathcal{M} \to \mathbb{R}^N$ & Grounding on base (well-defined by Axiom 4) \\
$\Sigma : \mathcal{P} \to V$ & Evaluation function \\
$\bar{\Sigma} : \mathcal{M} \to V$ & Evaluation on base (well-defined by Axiom 7) \\
$\mathcal{C} \subseteq \mathcal{M}$ & Admissible region (constraint set) \\
$\mathcal{U}$ & Control space \\
$u : \mathcal{M} \to \mathcal{U}$ & BIP-compliant control law \\
$L : T\mathcal{P} \to \mathbb{R}$ & Decision Lagrangian \\
$S[\gamma]$ & Action functional \\
$J \in \mathfrak{g}$ & Alignment current (Noether charge) \\
$\mathcal{S} = (\mathbb{R}_{>0})^d$ & Scaling group (unit changes) \\
$\Psi^*$ & Dimensionless grounding \\
\bottomrule
\end{tabular}
\end{table}

\begin{thebibliography}{9}

\bibitem{bond2026ge}
A.~H.~Bond, \emph{Geometric Ethics: The Mathematical Structure of Moral Reasoning}, v1.14, Department of Computer Engineering, San Jos\'e State University, Spring 2026.

\bibitem{bond2026erisml}
A.~H.~Bond, ``ErisML/DEME v3.0: A Library for Governed, Foundation-Model-Enabled Agents with Tensorial Ethics,'' \url{https://github.com/ahb-sjsu/erisml-lib}, 2026.

\end{thebibliography}

\end{document}
